\section{Proof of the main theorem}
\label{sec:main-proof}

We apply orbit cancellation in each conditional fibre and integrate the
measurable finite-block averages from Proposition~\ref{prop:block-witness}.

\begin{proof}[Proof of Theorem~\ref{thm:main}]
Suppose the conclusion fails, and let $k\ge3$ be the least failing order.
Proposition~\ref{prop:zero-entropy-witness} gives a mixing finite-alphabet
process $(X,\mu_X,T_X)$ of zero entropy, centered finite-valued one-site
functions, and a sequence of layouts witnessing failure at order $k$.
All proper sublayouts satisfy the lower-order mixing conclusions.

Construct the ordered array laws and their saturated extension by
Section~\ref{sec:arrays}. The square locality and operator results of
Sections~\ref{sec:arrays} and~\ref{sec:operators} apply to this extension.
Proposition~\ref{prop:block-witness} supplies a centered one-site witness
$f$, a countable abelian group $\mathcal L$, an automorphism $G$, a finite
nonempty set $B\subset\mathcal L$, and representatives $\chi_l^u$ in each
typical fibre. The measurable coefficient sums $s_r(u)$ defined in
\eqref{eq:block-witness-sum} satisfy
\[
 0\le s_r(u)\le |B|\norm f_\infty,
 \qquad \int_U s_r\dd\nu=c>0\qquad(r\ge0).
\]

Fix a typical $u$. The marginal of $x$ under $\eta_u$ is exactly $\mu_X$.
The representatives obey \eqref{eq:projective-characters} by
\eqref{eq:block-character-laws}. Apply
Proposition~\ref{prop:orbit-cancellation} to each of the finitely many
$l\in B$. It gives
\[
 \frac1N\sum_{r<N}s_r(u)\longrightarrow0.
\]
This is a pointwise argument in the fibre; it uses neither invariance of
$\eta_u$ under time nor a measurable choice of the abstract identification.
The original sums are already measurable, and are bounded by
$|B|\norm f_\infty$ for every $r$ and almost every $u$.
Bounded convergence therefore makes their integrals tend to zero.
Every such integrated average equals $c>0$, a contradiction.

There is consequently no least failing finite order. Since ordinary
mixing supplies order two, all the conclusions in
Theorem~\ref{thm:main} follow.
\end{proof}
